\documentclass[main]{subfiles}

\begin{document}
\ifSubfilesClassLoaded{\setcounter{chapter}{5}}{}
\chapter{Conclusion}

\section{Summary of the Thesis}

This thesis set out to close a specific gap identified in Chapter~1: the
Standard Model Extension (SME) and the Dobrescu--Mocioiu (DM) classification
of exotic spin-dependent potentials describe the same physics in two
languages, one relativistic and field-theoretic, the other
non-relativistic and directly tied to laboratory observables, with no
complete, systematic dictionary between them, and no framework for using
that dictionary to compare matter- and antimatter-sector constraints as a
test of CPT symmetry. Three intertwined pieces of work followed from this.
Chapters~2 and~3 built the theoretical bridge, deriving the
non-relativistic Hamiltonians generated by the SME coefficients $b_\mu$,
$H_{\mu\nu}$, and $d_{\mu\nu}$ via the Foldy--Wouthuysen (FW) transformation
and matching each onto the Dobrescu--Mocioiu potential catalogue
$V_1$--$V_{16}$. Chapter~4 applied that bridge to real data, compiling 273
experimental constraint datasets into a single standardised database and
computing the CPT asymmetry parameter $A_\alpha$ for every
matter--antimatter pair the database could support. Chapter~5 then turned
the same compiled database around to ask where it is currently silent:
which fermion sectors and which DM potentials have no antimatter-sector
constraint at all, and used that answer to propose concrete next steps,
both experimental and archival.

The result is not a single number or a single discovery, but a working
pipeline: a mapping table connecting relativistic coefficients to
measurable potentials, a standardised constraint database built to be
extended, and a quantitative account of exactly how far current data can
and cannot go toward testing CPT symmetry in the spin-dependent sector.
The remainder of this chapter revisits the five research objectives set out
in Chapter~1 against what was actually delivered, draws together the
principal findings and their limitations, and outlines the theoretical,
database, and experimental work that would extend this framework further.

\section{Revisiting the Research Objectives}

\begin{enumerate}
  \item \emph{Derive explicit mappings between SME coefficients
    $(b_\mu, d_{\mu\nu}, H_{\mu\nu})$ and DM coupling structures using the
    FW transformation.} Achieved for $b_\mu$ and $H_{\mu\nu}$ in full: both
    derivations reproduce the expected operator structure and CPT
    behaviour at every order in $1/m$ kept in this thesis, and are backed
    by executable symbolic computation. Achieved only partially for
    $d_{\mu\nu}$: the mass-enhanced $d_{i0}\to V_2$ result matches
    Kostelecký and Lane (1999) exactly, including its overall sign, and can
    be cited as an established result; the momentum-dependent $d_{ij}\to
    V_7,V_8$ and $d_{00}\to V_8$ results are correct in operator structure
    but carry an overall sign this thesis's equation-of-motion-level
    treatment does not resolve (Section~\ref{sec:limitations}).
  \item \emph{Construct a unified, standardised database of experimental
    constraints.} Achieved: 273 datasets spanning NV-centre magnetometry,
    torsion balances, atomic co-magnetometers, positronium spectroscopy,
    and antimatter experiments were compiled, unit-standardised to a common
    $\lambda$ axis, and classified by coupling family, DM potential, and
    fermion sector.
  \item \emph{Define and compute the CPT asymmetry parameter $A_\alpha$
    across all relevant channels.} Achieved, with an important
    qualification that is itself one of this thesis's principal findings:
    $A_\alpha$ was computed for all ten matter--antimatter pairs the
    database could support, but Chapter~3 shows analytically, and
    Chapter~4 confirms empirically, that $A_\alpha$ computed from one-sided
    experimental upper bounds cannot by itself distinguish genuine CPT
    violation from an ordinary sensitivity gap between the two experiments
    being compared (Section~\ref{sec:findings-asymmetry}).
  \item \emph{Perform a global parameter-space analysis identifying
    unconstrained or weakly constrained regions.} Achieved: Chapter~5
    quantifies the antimatter-sector coverage gap by both fermion sector
    and DM potential, identifying the $e$-$N$ sector and the
    $V_{1a}$, $V_{9+10}$, $V_1$, and $V_{12+13}$ potentials as the largest
    single opportunities.
  \item \emph{Propose optimised experimental strategies for future CPT and
    exotic-interaction searches.} Achieved: Chapter~5 proposes prioritising
    new antimatter-sector bounds in the identified gap sectors and
    potentials, and, following the precedent set by Cong et
    al.~\cite{Cong2025}, reinterpreting existing antimatter- and
    Lorentz/CPT-sensitive datasets before commissioning new measurements.
\end{enumerate}

Four of the five objectives were achieved in full; the fifth ($d_{\mu\nu}$
mapping) was achieved for its dominant, mass-enhanced component and left
honestly open for its two subleading components. This is consistent with
the standard this thesis has tried to hold throughout: reporting what the
derivation actually shows, rather than what would make a tidier final
table.

\section{Principal Findings}

\subsection{The SME--Dobrescu--Mocioiu Mapping}

The master matching table of Chapter~3 (Table~3.4) is the thesis's central
technical deliverable: it gives, for each of $b_i$, $b_0$, $H_{ij}$,
$H_{0i}$, $d_{i0}$, $d_{ij}$, and $d_{00}$, the resulting non-relativistic
Hamiltonian, the DM potential(s) it generates, its order in $1/m$, and its
CPT parity. The single most physically significant result in this table is
the $b_i$ derivation: a genuinely CPT-odd coefficient produces an exact
sign flip, $H_{NR}=-\mathbf b\cdot\boldsymbol\sigma$ for matter versus
$+\mathbf b\cdot\boldsymbol\sigma$ for antimatter, obtained independently
from the Lagrangian's CPT transformation and from an explicit
charge-conjugation calculation on the Dirac spinor. This is, as far as this
thesis has been able to establish, the only coefficient among those treated
here for which the matter--antimatter difference is a genuine, sign-level
prediction rather than an artefact of comparing bounds of unequal
precision, which makes the epistemological result of
Section~\ref{sec:findings-asymmetry} below all the more important, since it
shows that even this genuine prediction cannot currently be read off the
compiled data unambiguously.

\subsection{The Asymmetry Parameter and the Sensitivity-Gap Problem}
\label{sec:findings-asymmetry}

The most conceptually important result of this thesis is not a numerical
bound but a methodological one, established in Chapter~3 and confirmed
against every matched pair in Chapter~4: the asymmetry parameter
$A_\alpha=(g_f-g_{\bar f})/(g_f+g_{\bar f})$, as actually computed by
SPINDEP from published experimental limits, cannot distinguish a genuine
CPT-violating signal from an ordinary difference in sensitivity between the
matter- and antimatter-sector experiments being compared. Substituting an
exact, signed CPT-odd relation ($g_{\bar f}=-g_f$) into this formula does
not give $|A_\alpha|\to1$, as might naively be assumed: it makes the
denominator vanish identically, so the expression diverges rather than
saturating, a result confirmed with a symbolic limit calculation. What
actually drives real SPINDEP output to $|A_\alpha|\approx1$ is the ordinary
behaviour of a ratio of two independent, always-positive upper bounds of
very different tightness, a pattern that arises regardless of the true
CPT parity of the underlying physics. This was not a purely theoretical
concern: it is exactly what Chapter~4 found in the compiled database, where
the CPT-odd $g_Ag_A$ pairs and the CPT-even $g_sg_s$ pair alike showed
$|A_\alpha|$ values that are statistically significant at effectively
$p\to0$ after correcting for the correlated nature of interpolated
constraint points, yet none of which can currently be read as evidence of
CPT violation rather than sensitivity-gap saturation. The single most
informative data point in the entire compiled database, by this standard,
is the lowest asymmetry in Table~4.2, the $g_Ag_A\cdot V_2\cdot
e$-$e^+$ pair at $|A_\alpha|=0.334$, precisely because it is the one pair
whose two bounds happen to be of comparable sensitivity, and so is not
merely reproducing the generic sensitivity-gap pattern.

\subsection{The Compiled Constraint Database}

Of the 273 datasets compiled, only 19 involve an antimatter sector at all
(9 in $e$-$\bar p$, 5 in $e$-$e^+$, 3 in $e$-$\bar\mu$, and one each in
$\bar p$-He and $dd\mu^+$), and from these only ten matter--antimatter
pairs could be matched by SPINDEP's pairing logic.
Every one of the ten is statistically significant after the
autocorrelation correction developed in Chapter~4, with effective degrees
of freedom of 6--21, a factor of 14--50 below the naive 300 interpolation
points, confirming that the correction is not a negligible refinement.
Taken together with the sensitivity-gap result above, the database
compilation and the statistical analysis point to the same conclusion from
two different directions: the limiting factor in testing CPT symmetry via
exotic spin-dependent interactions is not statistical power, but the sheer
scarcity and unevenness of antimatter-sector data.

\subsection{The Antimatter-Sector Coverage Gap}

Chapter~5's sector- and potential-level breakdowns make this scarcity
concrete. Six of the eleven classified DM potentials,
$V_{11}$, $V_{15}$, $V_{16}$, $V_{1a}$, $V_{9+10}$, and $V_{12+13}$, have
\emph{zero} antimatter-sector datasets despite having 2--31 matter-sector
datasets each, and four fermion-sector pairs, $e$-$N$, $e$-$n$,
$n$-$p$, and $n$-$n$, together account for 91 matter-sector datasets and
not one antimatter-sector counterpart. The $e$-$N$ sector alone has 58
matter-sector datasets, more than any other sector in the database, against
zero in $e$-$\bar N$. This is the binding constraint identified by this
thesis: not a lack of matter-sector precision, but an almost complete
absence of comparably precise antimatter-sector measurements in exactly the
channels where matter-sector data is most abundant.

\section{Limitations}
\label{sec:limitations}

Three limitations are worth stating together, since each bears on how much
weight the results above can be given.

\textbf{The $d_{\mu\nu}$ sign ambiguity.} The $d_{ij}\to V_7,V_8$ and
$d_{00}\to V_8$ derivations agree with Kostelecký and Lane (1999) in
operator structure but not in overall sign, a discrepancy traced to a
wavefunction-renormalisation step in the kinetic sector that this thesis's
equation-of-motion-level treatment does not include. Until this is
resolved, no experimental target should be motivated from the sign of
these two components specifically, only from their structure and order in
$1/m$.

\textbf{Database data quality.} Roughly a seventh of the compiled database
(39 of 273 datasets) carries a filename that cannot be resolved to a
specific $V_n$ token and sits in a directory that does not name the
potential either, and is recorded as \texttt{UNKNOWN} rather than guessed
at. A further eleven datasets, all under a single directory whose name
states the potential unambiguously, are classified as $V_1$ from that
directory rather than the filename; two more are confirmed as $V_1$ by
checking their source publications directly rather than inferred from a
coupling label. Two additional datasets carry the coupling label
\texttt{Code-plot-matlab}, a parser artefact rather than a physical
category, and a small number of entries under $V_{2+3}$ and $V_{9+10}$
report physically implausible interaction ranges ($\lambda\gtrsim10^{26}$~m)
almost certainly caused by a unit-detection error. None of these affect
the ten matched pairs analysed in Chapter~4, but they should be resolved
before any broader coverage claim is drawn from the full 273-dataset
database.

\textbf{Scope of the SME sector treated.} This thesis derives the
non-relativistic limit of only three SME coefficients,
$b_\mu$, $H_{\mu\nu}$, $d_{\mu\nu}$, out of the eight listed in
Chapter~2's classification table. The coefficients $a_\mu$, $c_{\mu\nu}$,
$e_\mu$, $f_\mu$, and $g_{\lambda\mu\nu}$ are discussed qualitatively for
completeness but not derived with the same rigour, and are not represented
in the SPINDEP asymmetry analysis.

\section{Future Work}

\subsection{Theoretical}

The most direct theoretical extension is to complete the $d_{\mu\nu}$
derivation with a wavefunction-renormalisation treatment of the kinetic
sector, which should resolve the open sign in the $d_{ij}$ and $d_{00}$
components and allow them to be cited with the same confidence as
$d_{i0}$. Beyond this, the FW machinery developed in Chapter~2 is general
and applies unchanged to the remaining SME coefficients not treated here;
extending the explicit mapping to $a_\mu$, $c_{\mu\nu}$, $e_\mu$, $f_\mu$,
and $g_{\lambda\mu\nu}$ would complete the SME--DM dictionary for the full
minimal (and near-minimal) fermion sector, rather than the subset most
directly relevant to spin-dependent forces.

\subsection{Database and Methodology}

The remaining 39 \texttt{UNKNOWN} datasets and the two
\texttt{Code-plot-matlab} entries are a concrete, bounded cleanup task:
resolving them requires checking each source publication individually,
since neither their filenames nor their directory paths carry a
resolvable signal. The implausible-$\lambda$ entries under $V_{2+3}$ and
$V_{9+10}$ should be
individually re-verified against source publications before the database
is treated as final. Separately, the $e$-$\mu$/$e$-$\bar\mu$ sector is
populated on both sides (three datasets each) but was not matched into a
pair by SPINDEP's pairing logic; determining whether this is because the
two sides lack a shared coupling/potential/interaction-class label or
because their $\lambda$ ranges do not overlap is a small, well-defined
follow-up that could add an eleventh matched pair to the database without
any new experimental input.

\subsection{Experimental}

Chapter~5 already develops the experimental strategy this thesis
recommends in detail: prioritising new antimatter-sector bounds in the
$e$-$N$, $n$-$p$, $n$-$n$, and $e$-$n$ sectors and in the $V_{1a}$,
$V_{9+10}$, $V_1$, and $V_{12+13}$ potentials, and searching existing BASE,
ALPHA, and co-magnetometer publications for data that could be
reinterpreted within this framework before new measurements are
commissioned, following the precedent Cong et al.~\cite{Cong2025} document
for comagnetometer data originally collected for sidereal Lorentz-violation
searches. The standard this thesis applies throughout is that a new
antimatter-sector bound is useful for a CPT test specifically when it
approaches the precision of its matter-sector counterpart, not merely when
it exists. A loose first measurement in an empty sector would only
reproduce the same sensitivity-gap-dominated $|A_\alpha|\approx1$ pattern
documented throughout Chapter~4, without resolving the CPT question it is
meant to address.

\section{Closing Remarks}

The gap this thesis set out to address, a missing dictionary between the
SME and the Dobrescu--Mocioiu potentials, and no systematic way to use that
dictionary to compare matter- and antimatter-sector constraints, is now
substantially, though not completely, closed. The mapping table of
Chapter~3, the compiled 273-dataset database of Chapter~4, and the gap
analysis of Chapter~5 together give a reproducible, extensible pipeline
rather than a fixed set of numbers: the database is built to absorb new
datasets as they are published, and the pairing and asymmetry-analysis code
apply unchanged to any new antimatter-sector measurement that lands in one
of the gap channels identified here. In that sense, the most durable
contribution of this thesis may not be any single derived Hamiltonian or
compiled bound, but the demonstration, backed by an explicit divergence
calculation and confirmed across every pair in the compiled database,
that an asymmetry parameter built from one-sided experimental bounds is a
sensitivity-gap diagnostic before it is a CPT diagnostic, and that closing
the antimatter-sector coverage gaps mapped in Chapter~5 is a precondition
for that parameter to mean what it was designed to mean. The translation
tables and analysis code developed here are offered as a practical
resource for that purpose, and, it is hoped, as a small contribution toward
the kind of unified reporting the SME data tables~\cite{Kostelecky2011}
were built to provide, a resource that CERN's antimatter precision
programme~\cite{Smorra2017,Ahmadi2017} is, at present, the platform best
placed to help complete.

\ifSubfilesClassLoaded{
\begin{thebibliography}{99}

\bibitem{Cong2025}
L.~Cong \textit{et al.},
\textit{Spin-dependent exotic interactions},
Rev.\ Mod.\ Phys.\ \textbf{97}, 025005 (2025).

\bibitem{Kostelecky2011}
V.~A. Kostelek\'{y} and N.~Russell,
\textit{Data tables for Lorentz and CPT violation},
Rev.\ Mod.\ Phys.\ \textbf{83}, 11 (2011), updated annually.

\bibitem{Smorra2017}
C.~Smorra \textit{et al.}\ (BASE Collaboration),
\textit{A parts-per-billion measurement of the antiproton magnetic moment},
Nature \textbf{550}, 371 (2017).

\bibitem{Ahmadi2017}
M.~Ahmadi \textit{et al.}\ (ALPHA Collaboration),
\textit{Observation of the 1S--2S transition in trapped antihydrogen},
Nature \textbf{541}, 506 (2017).

\end{thebibliography}
}{}

\end{document}
